\clearpage
\phantomsection
\section*{Chương 1. Nguyên lý Praat từ vật lý đến toán (Boersma 1993)}
\addcontentsline{toc}{section}{Chương 1. Nguyên lý Praat từ vật lý đến toán (Boersma 1993)}
\pdfbookmark[1]{Chương 1. Nguyên lý Praat từ vật lý đến toán (Boersma 1993)}{chuong1}

\begin{learningobj}
Chương này lần theo đúng mười hai bước của thuật toán đo cao độ $F_0$ trong phần mềm Praat (Boersma 1993; mã nguồn C++ \texttt{Sound\_to\_Pitch.cpp} và \texttt{Pitch.cpp}, nhánh autocorrelation). Với mỗi bước, bạn sẽ: hiểu hiện tượng vật lý mà bước đó bảo vệ; dẫn công thức từ định nghĩa gốc; tính tay một ví dụ số nhỏ đã kiểm bằng tập lệnh \texttt{verify/verify\_chuong\_00\_01.py}; thấy ánh xạ silicon trên kit DE10-Standard (giữ, xấp xỉ hay xóa); và nắm mô hình Python khớp bit cùng ngân sách ALM/DSP/M10K/RTF. Các khái niệm đã giải thích ở Chương~0 (ADC, Nyquist, aliasing, framing, mainlobe/sidelobe, Hanning, tự tương quan, Wiener--Khinchin, butterfly, BFP, Q$m.n$, Avalon-ST, RTF,\ldots) chỉ được \emph{tham chiếu}, không định nghĩa lại.
\end{learningobj}

\subsection*{1.0. Bản đồ mười hai bước: giấy Boersma $\leftrightarrow$ C++ $\leftrightarrow$ silicon}

Hình~\ref{fig:ch01_bp_map} tóm tắt chuỗi đầy đủ Praat AC và ghi chú ngay trên bản đồ những bước mà phần cứng streaming trên Cyclone~V sẽ giữ nguyên, xấp xỉ hoặc xóa có chủ đích. Kỹ sư phải hiểu bước đã xóa \emph{trước khi} xóa --- đó là nguyên tắc xuyên suốt Chương~1.

\begin{figure}[H]
\centering
\begin{tikzpicture}[font=\scriptsize, >=Stealth,
  box/.style={draw, rounded corners=2pt, align=center, inner sep=2.5pt, minimum height=0.55cm}]
  \foreach \i/\lab/\y/\col in {
    1/{Bước 1: Soft upsample (FFT $\times$ mask $\times$ IFFT $\times 2$)}/0/blue!6,
    2/{Bước 2: Khung $L$ + hop $H$ từ $F_{\min}$}/-0.65/blue!6,
    3/{Bước 3: Trừ mean cục bộ (local mean)}/-1.30/teal!10,
    4/{Bước 4: Cửa sổ Hanning (hoặc Gaussian)}/-1.95/teal!10,
    5/{Bước 5: Tự tương quan $\mathcal{R}_x$ (FFT hoặc MAC)}/-2.60/teal!10,
    6/{Bước 6: Chuẩn hóa $r_a[\tau]=(\mathcal{R}_x/\mathcal{R}_x[0])\,C[\tau]$}/-3.25/teal!10,
    7/{Bước 7: Cực đại cục bộ + nội suy parabolic}/-3.90/teal!10,
    8/{Bước 8: Tinh chỉnh Sinc / Brent}/-4.55/orange!12,
    9/{Bước 9: Octave cost (phạt thiên vị quãng tám)}/-5.20/teal!10,
    10/{Bước 10: Silence + voicing + ứng viên UV}/-5.85/teal!10,
    11/{Bước 11: Octave-jump + V$\leftrightarrow$UV costs}/-6.50/orange!12,
    12/{Bước 12: Viterbi \texttt{Pitch\_pathFinder}}/-7.15/orange!12} {
    \node[box, fill=\col, minimum width=9.0cm] (s\i) at (0,\y) {\i. \lab};
  }
  \foreach \i in {1,...,11} {\pgfmathtruncatemacro{\j}{\i+1}\draw[->, thick] (s\i)--(s\j);}
  \node[box, fill=gray!8, minimum width=5.2cm, text width=5.0cm, align=left, anchor=north]
    at (8.0,-0.1)
    {\textbf{Nguồn gốc:}\\
     Boersma 1993 §§1--4;\\
     \texttt{Sound\_to\_Pitch.cpp} (AC);\\
     \texttt{Pitch\_pathFinder} trong \texttt{Pitch.cpp}.\\[3pt]
     \textbf{FPGA mặc định (DE10):}\\
     Giữ bước 3--7, 9--10 dưới dạng MAC + ROM + ngưỡng.\\
     Xóa bước 1 (soft upsample cả file).\\
     Xấp xỉ bước 8 (Sinc $\to$ parabolic).\\
     Xóa/đẩy HPS bước 11--12 (Viterbi).};
\end{tikzpicture}
\caption{Chuỗi mười hai bước Praat autocorrelation. Ô xanh lá: giữ gần bit-exact trên fabric. Ô cam: xấp xỉ hoặc đẩy sang HPS. Bước~1 (ô xanh dương trên cùng) bị xóa trên fabric và thay bằng FIR chống alias ở Chương~4.}
\label{fig:ch01_bp_map}
\end{figure}

\subsection*{1.1. Bước 1 --- Soft upsampling: dọn mép Nyquist trước khi cắt khung}

\begin{firstprin}{Hiện tượng thật: năng lượng sát Nyquist bị cửa sổ ``quăng'' xuống vùng $F_0$}
Chương~0 đã giải thích lấy mẫu, tần số Nyquist $f_N = f_s/2$, aliasing và rò phổ (mainlobe/sidelobe) của cửa sổ Hanning. Hệ quả thực tế đối với đo cao độ là: nếu một phát ngôn (\emph{utterance}) còn năng lượng gần $f_N$ (tiếng rít quạt, nhiễu cơ học, sóng hài cao), thì khi nhân cửa sổ Hanning dài $40\,\text{ms}$, các thùy phụ khoảng $-32\,\text{dB}$ sẽ phát tán năng lượng đó xuống đúng dải tìm kiếm $75\ldots 600\,\text{Hz}$. Thuật toán sẽ thấy các đỉnh tự tương quan giả --- không phải dây thanh tạo ra chúng.

Boersma (1993, Step~1) xử lý cả một utterance trên máy tính như sau: (1)~tính FFT toàn chuỗi; (2)~nhân phổ với mặt nạ mềm $G[k]$ giảm tuyến tính từ $0{,}95\cdot f_N$ xuống $0$ đúng tại $f_N$; (3)~chèn zero giữa nửa phổ dương và nửa phổ âm rồi IFFT dài $2N$. Kết quả là upsample ``mềm'' $\times 2$ đồng thời đã dọn sạch mép Nyquist. Đây là thao tác tính toán trên cả câu nói nằm trong RAM, không phải hiện tượng vật lý mới của dây thanh.
\end{firstprin}

\paragraph{Toán từng bước.}
Đặt $f_N = f_s/2$. Mặt nạ phổ được định nghĩa theo bin tần số $k$ (với $N$ điểm FFT, bin Nyquist là $k = N/2$):
\[
G[k] =
\begin{cases}
1, & \dfrac{|k|}{N} < 0{,}95\cdot\dfrac{1}{2},\\[6pt]
\dfrac{f_N - |k|\,f_s/N}{f_N - 0{,}95\,f_N}, & 0{,}95\,f_N \le |k|\,f_s/N \le f_N,\\[6pt]
0, & |k|\,f_s/N > f_N.
\end{cases}
\]
Chuỗi sau xử lý là $x_{\mathrm{up}} = \mathrm{IFFT}_{2N}\bigl\{\mathrm{pad}\bigl(G[k]\,X[k]\bigr)\bigr\}$. Phép zero-pad trong miền tần số rồi IFFT dài gấp đôi chính là nội suy Fourier lý tưởng: lưới thời gian dày gấp hai, phổ hữu ích dưới Nyquist cũ được giữ nguyên (xem lại định nghĩa zero-padding ở Chương~0).

\paragraph{Ví dụ tính tay phổ 8 điểm.}
Lấy tín hiệu ``cứng'' đúng tại tần số Nyquist của lưới 8 điểm: $x = [1,0,-1,0,1,0,-1,0]$. DFT 8 điểm (không chuẩn hóa) cho năng lượng tập trung ở bin gần mép (bin $k=2$ và đối xứng), còn bin Nyquist $k=4$ bằng 0 trên tín hiệu này vì chu kỳ khớp lưới:
\begin{center}\small
\begin{tabular}{@{}rrrr@{}}\toprule
$k$ & $\mathrm{Re}\,X[k]$ & $\mathrm{Im}\,X[k]$ & $|X[k]|^2$ \\
\midrule
0 & 0 & 0 & 0 \\
1 & 0 & 0 & 0 \\
2 & 4 & 0 & 16 \\
3 & 0 & 0 & 0 \\
4 (Nyquist) & 0 & 0 & 0 \\
\bottomrule
\end{tabular}
\end{center}
Soft filter nhân mỗi bin với $G[k]$: giữ nguyên khi $k/N < 0{,}95/2$, rồi giảm tuyến tính về 0 tại bin Nyquist. Sau đó IFFT dài 16 mẫu --- tương đương nội suy dày gấp đôi trên lưới thời gian. (Các số DFT trên đã kiểm bằng \texttt{numpy.fft.fft}; xem \texttt{verify/verify\_chuong\_00\_01.py}.)

\begin{figure}[H]
\centering
\begin{tikzpicture}[font=\small, >=Stealth]
  \draw[->] (0,0)--(7.2,0) node[right]{$f$};
  \draw[->] (0,0)--(0,2.4) node[above]{$G(f)$};
  \draw[very thick, teal!80!black] (0,1.9)--(4.5,1.9)--(5.6,0)--(6.5,0);
  \node[below] at (4.5,0) {$0{,}95\,f_N$};
  \node[below] at (5.6,0) {$f_N$};
  \node[teal!80!black] at (2.2,2.15) {giữ nguyên};
  \draw[<->, orange!80!black, thick] (4.5,0.15)--(5.6,0.15);
  \node[orange!80!black, font=\scriptsize] at (5.05,0.45) {dốc tuyến tính};
\end{tikzpicture}
\caption{Mặt nạ phổ mềm Boersma Step~1: giữ nguyên dưới $0{,}95\,f_N$, dốc tuyến tính về 0 tại Nyquist, rồi zero-pad và IFFT dài $2N$ để upsample $\times 2$.}
\label{fig:ch01_soft_mask}
\end{figure}

\paragraph{Ý nghĩa phần cứng.}
Trên FPGA streaming, mẫu âm thanh đến từng nhịp đồng hồ: không thể đợi hết cả câu nói rồi FFT toàn utterance. DE10 thay bước~1 bằng (1)~lọc chống alias analog trước ADC của codec WM8731, và (2)~FIR/decimation sạch từ $48\,\text{kHz}$ xuống $16\,\text{kHz}$ (Chương~4, FIR Compiler II). Cùng mục tiêu vật lý --- đừng để năng lượng gần Nyquist giả thành chu kỳ thấp trong vùng $F_0$ --- khác công cụ.

\paragraph{Twin Python bit-exact.}
Twin mặc định \emph{bỏ qua} soft upsample và bắt đầu từ PCM đã anti-alias. Lab đối chiếu Praat đầy đủ có thể chạy offline:
\texttt{np.fft.rfft} $\to$ nhân $G[k]$ (cùng dốc $0{,}95\,f_N\to f_N$) $\to$ \texttt{irfft(n=2N)}, rồi mới đưa vào đường pitch. Sai điểm cắt mask sẽ lệch năng lượng mép phổ so với Praat gốc.

\begin{costbox}
\textbf{Ngân sách silicon --- Bước 1 Soft upsample.}
\begin{itemize}[leftmargin=1.2em]
\item Trên fabric streaming: \textbf{0 ALM / 0 DSP / 0 M10K} cho Step~1 (không thực hiện FFT cả utterance).
\item Thay thế: anti-alias analog + FIR Compiler II (Chương~4) --- vài chục khối DSP và vài M10K hệ số lọc.
\item Twin đối chiếu Praat: soft-upsample offline trên Python/HPS trước khi đưa PCM vào đường fixed-point --- không tính vào RTF thời gian thực.
\item Quên dọn Nyquist: tone $\sim 7\ldots 8\,\text{kHz}$ giả thành chu kỳ dài trong vùng $75\ldots 600\,\text{Hz}$, làm sai $F_0$ dù khối MAC phía sau hoàn toàn đúng.
\end{itemize}
\end{costbox}

\paragraph{Bài tập 1.1.} Một utterance tại $f_s = 16\,\text{kHz}$ còn nhiễu $9\,\text{kHz}$. (a)~Tần số alias sau ADC là bao nhiêu? (b)~Soft upsample của Praat có ngăn được aliasing \emph{đã xảy ra} ở ADC không, hay chỉ dọn năng lượng còn lại dưới Nyquist trước khi cửa sổ hóa? Giải thích.

\paragraph{Lời giải.} (a)~$f_{\mathrm{alias}} = |16 - 9| = 7\,\text{kHz}$ (đã giải ở Chương~0). (b)~Soft upsample \emph{không} đảo ngược aliasing đã xảy ra trong ADC: thành phần $7\,\text{kHz}$ đã nằm trong dải gốc. Nó chỉ cắt/làm mềm phần phổ còn lại gần $f_N$ \emph{trước khi} nhân cửa sổ, giảm rò phổ xuống vùng $F_0$. Trên FPGA, việc chống alias thật sự phải làm bằng lọc analog + FIR trước hạ mẫu (Chương~4).

